\begin{table}[htbp]
  \centering
  \caption{Exact type I error rate as a function of the common response probability assumed under the null hypothesis, at $n_{1} = n_{0} = 192$ with $\pi_{1} = 0.45$, $\pi_{0} = 0.30$, $\omega_{1} = \omega_{0} = 0.10$, $\rho_{1} = \rho_{0} = 0$ and $\alpha = 0.025$ (one-sided). The size of a two-sample binomial test is not constant in this nuisance parameter, so no single value establishes size control.}
  \label{tab:null-prevalence}
  \begin{threeparttable}
  \begin{tabular}{llr}
    \toprule
    Analysis & Assumed null probability & Type I error rate \\
    \midrule
    NRI & 0.1000 & 0.0252 \\
     & 0.2000 & 0.0249 \\
     & 0.2700 & 0.0248 \\
     & 0.3000 & 0.0251 \\
     & 0.3375 (weighted average) & 0.0248 \\
     & 0.4050 & 0.0260 \\
     & 0.5000 & 0.0232 \\
    \addlinespace
    CC & 0.1000 & 0.0250 \\
     & 0.2000 & 0.0252 \\
     & 0.3000 & 0.0251 \\
     & 0.3750 (weighted average) & 0.0251 \\
     & 0.4500 & 0.0251 \\
     & 0.5000 & 0.0251 \\
    \bottomrule
  \end{tabular}
  \begin{tablenotes}[flushleft]\footnotesize
    \item The sample size is the one the NRI formula gives under these assumptions. For the NRI analysis the assumed probability is on the NRI scale; for the CC~analysis it is the conditional response probability among completers. The row marked as the weighted average is the convention used in the main text.
  \end{tablenotes}
  \end{threeparttable}
\end{table}
